\chapter{実験と考察}
\label{ch:experiments}
\label{chap:experiments}

本章では、Eppstein アルゴリズムの最外層並列化を実装し、その性能を測定した過程をたどる。
話の順序は、実際に測定を進めた時系列にそろえてある。
ひとつの測定が次の問いを生み、その問いに答えるために次の測定を設計する、という連鎖として書く。
逐次の二アルゴリズムの比較から始め、素朴な並列化がなぜ頭打ちになるかを負荷分布とフェーズ別の内訳で突き止め、配分戦略・非対称コア・共有メモリ化と順に切り込んでいく。

各節は、何を知りたいか（意図）、どう測ったか（方法）、何が出たか（結果）、それが何を意味するか（考察）の順で述べる。
測定値には、直接測ったものと、第三者の測定や単純なモデルから導いた推定とがある。
両者は文中で区別し、推定にはその根拠を添える。

\section{実験環境と測定プロトコル}
\label{sec:exp-env}

\subsection{意図}

並列性能の測定は環境に敏感である。
本研究のように単一のノート型計算機で測る場合、背景で動くアプリケーションや熱の状態が数値を大きく揺らす。
どの数値がアルゴリズムの性質を表し、どの数値が測定環境の副作用かを後から切り分けられるよう、測定条件と手続きをここで固定しておく。

\subsection{実験環境}

測定はすべて Apple M4 を搭載した MacBook 1 台で行った。
M4 は高速な Performance コア（以下 P コア）4 個と低速な Efficiency コア（以下 E コア）6 個の非対称構成で、主記憶は 16\,GB である。
実装は外部の数値ライブラリを用いない純 Python で、並列化には標準ライブラリの \texttt{multiprocessing} を使う。
P コアと E コアの速度差、および 16\,GB という主記憶の制約は、後の節（\ref{sec:pe-split}、\ref{sec:ecore-speed}）で性能の説明そのものに関わってくる。

\subsection{測定プロトコル}

計測値の揺れに対しては、次の三つの手続きで対処する。

\begin{itemize}
  \item \textbf{反復最小値}：同一条件を 3 回繰り返し、最小値を採用する。実行時間は他プロセスの割り込みで上振れする一方、下振れはしない。最小値は割り込みの影響が最も小さい実行に対応する。
  \item \textbf{環境スナップショット}：測定の直前に \texttt{ps} で背景プロセスの状態を記録し、想定外の負荷が走っていないことを確認してから実行する。
  \item \textbf{時間的な隣接}：比較したい条件どうしは、同一セッション内で時間を空けずに連続して測る。後述する熱ドリフトが条件間に混入するのを防ぐためである。
\end{itemize}

このプロトコルが必要なことは、環境の影響を定量した二つの観測が裏づける。
ひとつは背景アプリの影響で、Chrome を閉じただけで soc-sign-epinions の並列列挙時間が約 26\% 短縮した（\ref{sec:pe-split} 節で詳述）。
もうひとつは熱ドリフトで、約 1.5 時間の連続測定のあと、同一の純計算ループが 0.64 秒から 0.85 秒へ 32\% 遅くなった（\ref{sec:ecore-speed} 節）。
どちらもアルゴリズムとは無関係な要因であり、この幅を条件間の差と取り違えないための手続きが上の三点である。

\section{逐次版の比較：Tomita 対 Eppstein}
\label{sec:seq-compare}

\subsection{意図}

並列化の土台として、まず逐次の実装を比較する。
本研究は疎グラフに有利な Eppstein を並列化の対象に選ぶが、その選択が妥当かを 10 種類のグラフで確かめておく。
比較の軸は\textbf{縮退数} $d$ である。
$d$ は、どの頂点も縮退順序で自分より後ろの近傍が $d$ 個以下になるように頂点を並べられる最小値で、グラフの疎さを表す。
Eppstein は最外層をこの縮退順序で回ることで計算量を $O(d\,n\,3^{d/3})$ に抑える。

\subsection{方法}

比較対象は、Tomita らのピボット版 CLIQUES と Eppstein の 2 実装である。
両者を同一プロセス内で交互に 3 回ずつ実行し、それぞれの最小値を採る。
交互反復にするのは、このマシンでは単発計測が 2〜3 割振れることがあるためで、交互に測れば環境の変動が両者に等しく乗り、相対比較が安定する。
3 回の振れは全データセットで 1\% 未満だった。
グラフは、生成規則が既知の合成グラフ 7 種と、Stanford の SNAP が配布する実データ 3 種からなる。

\subsection{結果}

結果を表~\ref{tab:seq-compare} に示す。
倍率は Tomita の時間を Eppstein の時間で割った値で、1 を超えれば Eppstein が速い。

\begin{table}[htbp]
  \centering
  \caption{逐次版の比較（各値は 3 回の最小値、倍率は Tomita/Eppstein）}
  \label{tab:seq-compare}
  \begin{tabular}{lrrrrrr}
    \toprule
    グラフ & $n$ & $d$ & 極大クリーク数 & Tomita & Eppstein & 倍率 \\
    \midrule
    木                       & 20{,}000  & 1   & 20.0千   & 0.502\,s & 0.030\,s & 16.6 \\
    2 次元格子 $100\times100$ & 10{,}000  & 2   & 19.8千   & 0.137\,s & 0.022\,s & 6.2 \\
    疎 Erdős–Rényi（平均次数10）& 20{,}000 & 7   & 99.7千   & 0.621\,s & 0.142\,s & 4.4 \\
    スケールフリー BA（$m{=}5$） & 20{,}000 & 5   & 96.9千   & 0.717\,s & 0.135\,s & 5.3 \\
    ca-AstroPh               & 18{,}772  & 56  & 36.4千   & 2.290\,s & 0.871\,s & 2.6 \\
    ca-HepPh                 & 12{,}008  & 238 & 14.9千   & 1.008\,s & 0.636\,s & 1.6 \\
    soc-sign-epinions        & 131{,}828 & 121 & 2{,}223万 & 167.3\,s & 94.9\,s & 1.8 \\
    密 Erdős–Rényi（$n{=}120$）& 120      & 48  & 3.5万    & 0.123\,s & 0.109\,s & 1.1 \\
    Moon–Moser（$k{=}11$）    & 33        & 30  & 17.7万   & 0.147\,s & 0.120\,s & 1.2 \\
    完全グラフ $K_{300}$      & 300       & 299 & 1       & 0.074\,s & 0.192\,s & 0.4 \\
    \bottomrule
  \end{tabular}
\end{table}

\subsection{考察}

結果は $d/n$（縮退数がグラフサイズに占める割合）でほぼ説明できる。
疎で大きいグラフほど Eppstein が勝ち、密になると互角に収束し、明確に負けるのは完全グラフだけである。
木（$d{=}1$）では 16.6 倍、完全グラフ（$d{=}299$）では 0.4 倍と、勝敗が $d$ の小ささにそのまま従っている。
Eppstein が最外層を縮退順序で回って最外層を小さな部分問題に分割できるのは、$d$ が小さい疎グラフに限られ、密グラフでは分割の余地がないためこの利点が消える。

以降の並列化はすべて Eppstein を対象とする。
主力のベンチマークには soc-sign-epinions を選ぶ。
逐次で約 94.8 秒かかり、並列化の効果を測るのに十分重く、かつ実データである点が理由である。

\section{素朴な並列化とバッチサイズ}
\label{sec:naive}

\subsection{意図}

最外層の各頂点を根とするサブ問題は互いに独立に処理できる。
そこで最外層で素朴に並列化し、worker 数を増やしたときにどこまで速くなるかを測る。

\subsection{方法}

グラフはプール起動時に各 worker へ一度だけ \texttt{pickle} 転送し、負荷分散は \texttt{imap\_unordered} の動的配分に任せる。
サブ問題はバッチにまとめて配る。
まずバッチサイズ 512 で worker 数を 1, 2, 4, 8 と変えて測った。

\subsection{結果}

結果を表~\ref{tab:naive-batch} に示す。

\begin{table}[htbp]
  \centering
  \caption{soc-sign-epinions の素朴な並列化（列挙時間、workers=1 比）}
  \label{tab:naive-batch}
  \begin{tabular}{lrr}
    \toprule
    構成 & 列挙時間 & workers=1 比 \\
    \midrule
    workers=1（基準）        & 99.8\,s & 1.00 \\
    workers=2、バッチ512     & 69.3\,s & 1.44 \\
    workers=4、バッチ512     & 67.1\,s & 1.49 \\
    workers=8、バッチ512     & 70.0\,s & 1.43 \\
    workers=8、バッチ64      & 27.7\,s & 3.60 \\
    \bottomrule
  \end{tabular}
\end{table}

バッチ 512 では 4 worker 以降がまったく伸びず、67〜70 秒で頭打ちになった。
変更したのはバッチサイズを 512 から 64 に細かくしたことだけで、それにより 8 worker が 70 秒から 27.7 秒まで縮んだ。

並列化が逆効果になる場合も確認した。
1 秒未満で終わるグラフ（例：木 $n{=}20{,}000$）では、プール起動とグラフの \texttt{pickle} 転送（合計約 0.3 秒）が支配的になり、すべてのケースで逐次より遅くなった。
ca-HepPh は固定費に加えて、239 頂点のクリークを含む探索木が 1 タスクに載るため、8 worker でも改善しなかった。

バッチ 512 の頭打ちは、負荷の偏りを示唆する。
そこで、頭打ちの理由を負荷分布の測定で確かめる。

\section{負荷分布の測定}
\label{sec:load-dist}

\subsection{意図}

バッチ 512 で並列化が頭打ちになった理由を明らかにする。
外側頂点ごとのサブ問題の所要時間を、全頂点で個別に計測した。

\subsection{結果}

結果を表~\ref{tab:load-dist} に示す。
数値は全所要時間に占めるシェアである。

\begin{table}[htbp]
  \centering
  \caption{サブ問題の負荷分布（全所要時間に占めるシェア）}
  \label{tab:load-dist}
  \begin{tabular}{lrrrrr}
    \toprule
    データセット & 頂点数 & 最重の1頂点 & 上位0.1\% & 上位1\% & 上位10\% \\
    \midrule
    soc-sign-epinions & 131{,}828 & 0.86\% & 45.2\% & 94.7\% & 99.7\% \\
    ca-HepPh          & 12{,}008  & 13.0\% & 20.0\% & 47.0\% & 90.1\% \\
    ca-AstroPh        & 18{,}772  & 0.21\% & 2.8\%  & 17.7\% & 65.9\% \\
    \bottomrule
  \end{tabular}
\end{table}

\subsection{考察}

soc では上位 1\% の頂点が仕事の 94.7\% を占め、その重い頂点は縮退順序の末尾に密集していた。
疎な外周から頂点を剥がしていくと密なコアが必然的に末尾に残るためである。
ただし最重の 1 頂点でも 0.86\% なので、バッチを細かくすれば負荷は均せる。
バッチ 512 から 64 への変更で 8 worker が 70 秒から 27.7 秒に縮んだのは、この偏りの解消にあたる。

ca-HepPh は性質が異なる。
1 頂点（239 頂点クリークの根）だけで全時間の 13.0\% を占め、この 1 タスクは分割しない限りどう配っても待たされる。
このためスピードアップの上限が約 7.7 倍に制限される。
「配り方で救える soc」と「サブ問題内部の分割が要る ca-HepPh」は、同じ「並列化が効かない」でも別の問題である。
バッチの細分化で救えるのは前者だけで、後者はタスクの粒度そのものが下限を作る。

\section{フェーズ別時間の分解}
\label{sec:phase}

\subsection{意図}

並列実行の時間がどこに消えているかを知る。
1 回の実行を 4 区間に分けて計測した。

\begin{itemize}
  \item \textbf{load}：グラフの読込・生成。並列実行の外側で 1 回だけ測る、worker 数に依存しない値。
  \item \textbf{ordering}：縮退順序の計算。逐次に行う準備ループ。
  \item \textbf{startup}：プロセスプールの生成と worker の初期化。
  \item \textbf{compute}：並列本体。
\end{itemize}

あわせて、同じ実行から worker 別の busy 時間と idle 率を取得した。
計測は 9 グラフ $\times$ workers=1,2,4,8 $\times$ 3 回（最小値採用）である。

\subsection{準備ループは軽い}

最初の問いは、逐次の準備ループ（縮退順序）がどれほど重いか、である。
答えは、心配するほど重くない、だった。
ordering が total に占める割合は全グラフ・全 worker 数を通して最大 17.8\%（絶対値 0.05 秒未満）で、主力の soc では 1\% 未満である。
Amdahl の逐次項としては当面無視してよい。

\subsection{固定費は startup に集中する}

かわりに固定費として効いていたのは startup だった。
プール起動時にグラフ全体を worker ごとに \texttt{pickle} 転送する設計の代償で、worker 数にほぼ比例して伸びる。
soc では w=1 の 0.27 秒から w=8 の 1.27 秒まで増える。
compute が 1 秒未満の 7 グラフでは、この startup のために w=8 が w=1 より遅くなる。

soc の worker 数スイープをフェーズ別に分けた内訳を表~\ref{tab:phase-soc} に示す。

\begin{table}[htbp]
  \centering
  \caption{soc-sign-epinions のフェーズ別内訳（バッチ64、3 回の最小値）}
  \label{tab:phase-soc}
  \begin{tabular}{rrrrrrrl}
    \toprule
    workers & ordering & startup & compute & total & 対w=1倍率 & idle\% & busy 最大/最小 \\
    \midrule
    1 & 0.432\,s & 0.268\,s & 126.006\,s & 126.706\,s & 1.00 & 0.0\%  & 126.0 / 126.0\,s \\
    2 & 0.366\,s & 0.338\,s & 65.800\,s  & 66.504\,s  & 1.91 & 6.9\%  & 65.7 / 56.8\,s \\
    4 & 0.339\,s & 0.643\,s & 48.419\,s  & 49.401\,s  & 2.56 & 15.0\% & 48.4 / 35.6\,s \\
    8 & 0.338\,s & 1.273\,s & 36.193\,s  & 37.803\,s  & 3.35 & 30.4\% & 36.2 / 11.8\,s \\
    \bottomrule
  \end{tabular}
\end{table}

startup が worker 数にほぼ比例して伸びているのが直接読める。
startup の値は近似である。
プール生成の直後に全 worker へ noop を配って初期化完了を待つ方式で測っているため、厳密な転送完了時刻とは一致しない。

\subsection{並列化が得なグラフは限られる}

9 グラフを通して見ると、並列化で得をしたのは逐次 compute が長い 3 グラフだけだった。
soc は worker を増やすほど伸びる（w=8 で 3.35 倍）のに対し、ca-AstroPh と ca-HepPh は w=4 がピーク（1.62 倍、1.51 倍）で、w=8 では startup の伸びが compute の削減を上回って逆に悪化する。
残りの 6 グラフ（compute が 0.2 秒以下）はすべての worker 数で損で、最悪の 2 次元格子は w=8 で w=1 の 5 倍遅い。

\subsection{スケールしない主因は compute 内の idle}

soc に話を戻すと、w=8 でも ordering と startup の合計は total の 4\% に過ぎない。
スケールしない主因は固定費ではなく、compute 内部の idle である。
idle\% は w=2 の 6.9\% から w=8 の 30.4\% へ伸び、busy 時間は最も忙しい worker の 36.2 秒に対し最も暇な worker は 11.8 秒だった。
負荷分布の測定で見つけた「末尾の重い頂点への偏り」が、worker の遊休として現れている。

次の問いは、この compute 内の idle を配り方だけでどこまで消せるか、である。

\section{配分戦略の比較}
\label{sec:alloc}

\subsection{意図}

フェーズ別分解で見つけた compute 内の遊休（w=8 で idle 30.4\%、busy 最大 36.2 秒に対し最小 11.8 秒）を、タスクの配り方だけでどこまで解消できるかを検証する。
先行研究の定石である work stealing や共有メモリ化は実装コストが大きい。
その前に、静的な配分の工夫で届く範囲を確かめておく位置づけである。

\subsection{方法}

比較した三つの戦略は、いずれも「縮退順序の各頂点を根とするサブ問題を、バッチにまとめて worker に配る」枠組みは同じで、バッチの作り方だけが違う。
前提として、負荷分布の測定から重い頂点は縮退順序の末尾に密集することがわかっている。

\begin{itemize}
  \item \textbf{block}（現行）：順序の先頭から連続にバッチサイズずつ区切る。重いバッチが最後に配られる。
  \item \textbf{reversed}（末尾逆順）：順序を逆にしてから連続に区切る。重いバッチを最初に配り、終盤の待ちを防ぐ。
  \item \textbf{interleave}（ストライド）：バッチ数を $K$ として、バッチ $k$ に頂点 $k, k+K, k+2K, \dots$ を入れる。どのバッチにも軽い頂点と重い頂点が混ざり、バッチの重さ自体がほぼ均一になる。
\end{itemize}

測定条件は soc-sign-epinions、バッチ 64、workers=4,8、各条件 3 回の最小値である。
単発計測の振れ対策として、block も同一セッションで再測定した基準を使う。
極大クリーク数は全条件で 22{,}226{,}173 に一致し、結果の正しさは確認済みである。

\subsection{結果}

結果を表~\ref{tab:alloc} に示す。

\begin{table}[htbp]
  \centering
  \caption{配分戦略の比較（soc-sign-epinions、バッチ64、3 回の最小値）}
  \label{tab:alloc}
  \begin{tabular}{llrrrr}
    \toprule
    戦略 & workers & compute & idle\% & busy 最大/最小 & busy 合計 \\
    \midrule
    block      & 4 & 33.3\,s & 20.2\% & 33.3 / 22.2\,s & 106\,s \\
    reversed   & 4 & 32.7\,s & 0.1\%  & 32.7 / 32.7\,s & 131\,s \\
    interleave & 4 & 36.7\,s & 0.1\%  & 36.7 / 36.6\,s & 147\,s \\
    block      & 8 & 25.9\,s & 32.6\% & 25.8 / 8.0\,s  & 139\,s \\
    reversed   & 8 & 28.5\,s & 23.5\% & 28.5 / 18.1\,s & 175\,s \\
    interleave & 8 & 24.8\,s & 0.1\%  & 24.7 / 24.6\,s & 198\,s \\
    \bottomrule
  \end{tabular}
\end{table}

w=8 の最速は interleave の 24.8 秒で、block からの短縮は 4.3\% にとどまった。
reversed は w=4 では僅かに勝つが、w=8 では block より悪化した。

\subsection{負荷の均等化は成功し、高速化は失敗した}

interleave は idle を 32.6\% から 0.1\% までほぼ完全に消した。
idle が純粋な無駄なら、compute は $1/(1-0.326)$ で約 1.5 倍縮むはずだったが、実際の短縮は 4.3\% だった。
理由は表の busy 合計に表れている。
idle が消えるのと引き換えに、worker が働いた時間の合計が 139 秒から 198 秒へ膨張しており、「遊んでいた worker を働かせる」効果を「1 単位の仕事にかかる時間が伸びる」効果がほぼ打ち消している。

膨張の有力な説明は、測定マシンが P コア 4 個と E コア 6 個の非対称構成であることだ。
w=8 では必ず E コアにも worker が載るため、worker の稼働時間は仕事量の物差しにならない。
idle\% という診断指標は「コアは均質である」という前提を暗黙に置いており、その前提が崩れるこのマシンでは、負荷均等化の効果を過大に予測していた。
もうひとつの候補として、重いサブ問題どうしを同時に走らせることによるキャッシュとメモリ帯域の競合も busy 膨張に寄与しうるが、こちらの切り分けは未検証である。

reversed が w=8 で悪化した機構は別である。
最重バッチを最初に引いた worker の busy（28.5 秒）が compute 全体と一致しており、最重バッチ 1 個がクリティカルパスになっている。
重い方から配る戦略は、バッチをさらに細かくしない限りこの下限を破れない。

帰結として、静的配分の改善はこのマシンではほぼ頭打ちである。
idle をほぼ 0\% にしてもなお高速化しなかったという事実は、次の二つの主張を検証する動機になる。
第一に、E コア 6 個で走る分の処理が全体の足を引っ張っているのではないか。
第二に、主記憶 16\,GB が不足し、swap や競合が起きているのではないか。

\section{P/E コアの切り分け実験}
\label{sec:pe-split}

\subsection{意図}

前節の二つの主張を検証する。
主張 1（P/E 非対称コアが足を引っ張っている）は、macOS の \texttt{taskpolicy -b} でプロセスを E コアに固定して同じ計算を走らせれば、E コアの速度を直接測れる。
主張 2（16\,GB 不足による swap・競合）は、グラフのメモリフットプリントの実測で swap の有無を判定できる。
あわせて、Chrome を閉じた環境で interleave の w=4,8 を再測定し、背景アプリの影響も確かめる。

\subsection{方法}

条件は soc-sign-epinions、interleave、バッチ 64、各条件 3 回の最小値である。
すべて同一セッションで、Chrome を閉じ、測定前に \texttt{ps} で背景プロセスがアイドルであることを確認してから実行した。
E コア条件は \texttt{taskpolicy -b} 経由で起動する。
background QoS に落とすことで Apple Silicon では E コアに固定されるが、この方法は優先度も同時に下げるため、測った E コア速度は下限値として扱う。
逆に P コア側を強制する手段は macOS になく、「重い単一プロセスは P コアに載る」というスケジューラの性質に依存している。

\subsection{結果}

結果を表~\ref{tab:pe-split} に示す。
スループットは w=1 通常の compute（96.0 秒）を当該条件の compute で割った値で、「P コア何個分の仕事が出たか」を表す。
極大クリーク数は全条件で一致する。

\begin{table}[htbp]
  \centering
  \caption{P/E コアの切り分け（soc-sign-epinions、interleave、3 回の最小値）}
  \label{tab:pe-split}
  \begin{tabular}{lrrrr}
    \toprule
    条件 & compute & スループット（P換算） & worker 1個あたり & 単独比の効率 \\
    \midrule
    w=1 通常（P）           & 96.0\,s  & 1.00 & 1.00 & ― \\
    w=1 E コア固定          & 363.6\,s & 0.26 & 0.26 & ― \\
    w=4 通常（実質 P のみ）  & 30.4\,s  & 3.16 & 0.79 & 79\% \\
    w=4 E コア固定          & 124.5\,s & 0.77 & 0.19 & 73\% \\
    w=8 通常（P+E 混載）     & 18.2\,s  & 5.26 & 0.66 & ― \\
    \bottomrule
  \end{tabular}
\end{table}

\subsection{主張 1 は実証された}

同じ計算が P コアで 96 秒、E コア固定で 363.6 秒と約 4 倍違う（速度比 0.26、ただし background QoS の影響を含む下限値）。
worker の稼働時間が仕事量の物差しにならないこと、すなわち配分戦略の比較で見た busy 膨張の主因が非対称コアであることが、推測ではなく測定値で言えるようになった。

\subsection{主張 2 は二つに分けて判定する}

swap 説は棄却できる。
グラフ 1 コピーのメモリ使用量を実測すると約 151\,MB で、worker 8 個と親を合わせても約 1.4\,GB と 16\,GB に十分収まる。

一方で共有資源の競合自体は実在する。
P コアだけの w=4 でも効率が 79\% に落ち、E コアだけの w=4 でも 73\% に落ちる。
idle はどちらもほぼ 0\% なので、これは待ちではなく「同じ種類のコアを複数同時に使うと 1 コアあたりの速度が下がる」現象である。
候補はメモリ帯域・キャッシュの共有と、複数コア稼働時のクロック引き下げで、この二つの内訳は未検証である。

\subsection{環境の影響も定量できた}

Chrome を閉じただけで interleave w=8 の compute は 24.8 秒から 18.2 秒へ約 26\% 縮んだ。
配分戦略の比較は全条件を同じ環境で測っているので相対比較の結論は揺るがないが、絶対値の比較は同一セッションでも背景アプリの状態まで揃える必要がある。
測定前に \texttt{ps} で環境を記録する手順をプロトコルに加えたのはこの観測による。

クリーンな環境での interleave w=8 は compute 18.2 秒（対 w=1 で 5.3 倍、total でも 4.9 倍）となり、このマシンでの到達点が更新された。
P コア換算 5.26 個分という値は、E コアの下限速度比 0.26 から作る単純モデル（$4 + 4\times0.26 \approx 5.1$）と整合する。
ただしこの一致は次節で見直す。

\section{E コアの速度の正体}
\label{sec:ecore-speed}

\subsection{意図}

切り分け実験で得た速度比 0.26 は、\texttt{taskpolicy -b} が優先度も同時に下げるため、「約 4 倍遅い」が E コアの実力なのかクロック制限の結果なのか区別できない。
この疑問（クロック比が機械的に出ているだけではないか）を検証する。

\subsection{純計算ループによる切り分け}

メモリをほぼ使わない整数演算ループを同じ方法で比較すると（3 回の最小値）、P コア 0.64 秒に対し E コア固定で 3.68 秒、速度比 5.75 倍だった。
soc の実測（3.79 倍）より差が大きい。
クリーク列挙にはメモリ待ちが含まれ、メモリ待ちの時間はクロックを落としても伸びないため、低クロックの被害が薄まる。
つまりこの差は「クロック制限が主因、メモリ待ちが緩和要因」という構図を支持する。
副産物として、この対比はクリーク列挙が計算律速寄りであることの傍証にもなっている。

\subsection{公表情報との突き合わせ}

Apple はスペックページにクロック周波数を載せない方針のため、公式値は存在しない。
代わりに、powermetrics による第三者測定（The Eclectic Light Company、Oakley 2024）がある。
M4 の E コアは 1{,}020〜2{,}592\,MHz の 7 段階 DVFS で動作し、background QoS のタスクは約 1{,}050\,MHz（最低周波数付近）に張り付くと報告されている。

これで実測値が分解できる。
P コア約 4.4\,GHz 対 1.05\,GHz のクロック比は 4.2 倍で、純計算ループの実測 5.75 倍との残差（$\times1.37$）が、E コアの IPC（1 クロックあたりの仕事量）が P コアより低い分と読める。
すると全力（2.59\,GHz）の E コアの速度比は $0.59 \div 1.37 \approx 0.43$ と推定できる。
これは直接測定ではなく、第三者測定のクロック値と自前の IPC 係数からの推定である。

\subsection{推定値でモデルを組み直す}

この $r \approx 0.43$ で w=8 をモデル化し直すと、理想のスループットは $4 + 4\times0.43 \approx 5.7$ P コア分になる。
実測は 5.26 なので効率 92\%、残り約 8\% が共有資源の競合ロスという内訳になる。
前節で 0.26 を使ったモデル（$4 + 4\times0.26 \approx 5.1$）が実測と合ったのは、E コアを遅く見積もる誤差と競合ロスを無視する誤差が打ち消し合った偽の一致だった。
いずれのモデルでも、配り方の改善ではこれ以上縮まないという結論は変わらない。

\subsection{方法論上の発見}

測定の過程で、E コア速度の測り方と熱ドリフトについて二つの知見を得た。

\begin{itemize}
  \item macOS ではユーザー空間の工夫でプロセスを「全力の E コア」に固定する手段が事実上ない。\texttt{taskpolicy -c utility} は P コアが空いていると E に固定されず（P と同速）、P コアを飽和させて押し出す方法もスケジューラがプロセスをコア間で回すため固定できない（1.19 倍しか落ちない）。確定させるには \texttt{sudo} で powermetrics を使い実クロックを観測するしかない（未実施、将来課題）。
  \item 長時間のベンチマーク連続実行でマシン自体が遅くなる。約 1.5 時間の測定後、同じソロの純計算ループが 0.64 秒から 0.85 秒へ 32\% 劣化した（熱によるクロック低下と見られるが pmset に記録はなく未確認）。同一セッション内でも時間が離れた測定の絶対値比較にはこのドリフトが乗る。比較したい条件は時間的に隣接させて測ることをプロトコルに加えた。
\end{itemize}

\section{共有メモリアプローチ}
\label{sec:shm}

\subsection{意図}

フェーズ別分解で、並列化の固定費が startup（グラフの \texttt{pickle} 転送）に集中することがわかった。
この固定費自体を消せば、フォーク済みの 6 グラフのように並列化が損だったケースを得に転じられる。
起動方式を変える二つのアプローチを試す。

\begin{itemize}
  \item \textbf{fork}：\texttt{multiprocessing.get\_context("fork")} で worker を起動する。子プロセスは親のメモリを copy-on-write で継承するため、グラフの \texttt{pickle} 転送も worker 側の再構築も発生しない。バッチの作り方は block と同一にし、block との差がすべて起動方式に帰着するようにした。macOS の既定が spawn なのは、スレッドを持つプロセスを fork すると子がデッドロックしうるためだが、本実装は fork 時点で単一スレッドの純 Python なので安全である。
  \item \textbf{shm}：起動方式は spawn のまま、グラフを CSR 形式の int64 配列 4 本（頂点 ID・indptr・隣接リスト・縮退順序）として \texttt{multiprocessing.shared\_memory} に置き、worker が名前で attach する。再帰カーネルは Python の set 演算に依存し、set は共有メモリに置けないため、worker は共有配列から自分用の dict-of-sets を再構築する。消えるのは \texttt{pickle} 化とパイプ転送だけで、worker 側の構造構築は残る。
\end{itemize}

spawn の startup を「interpreter 起動＋転送＋構築」に分解したとき、転送だけを消すのが shm、三つとも消すのが fork という対比になる。

\subsection{方法}

測定条件はバッチ 64、workers=1,4,8、各条件 3 回の最小値である。
同一セッションで fork、shm の順に連続実行した。
測定前に \texttt{ps} を確認し、前セッションの残骸である busy loop 2 プロセス（background QoS で E コアを占有）を停止してから開始した。
極大クリーク数は soc の 22{,}226{,}173 ほか全条件で一致する。
倍率は同一手法内の対 w=1 の total 比である。

\subsection{結果}

9 グラフの全数フェーズ別内訳は付録に回し、ここでは各グラフの w=4,8 での倍率だけを表~\ref{tab:shm-summary} にまとめる。
倍率が 1 を超えれば並列化の得である。

\todo{付録に全数表}

\begin{table}[htbp]
  \centering
  \caption{fork と shm の 9 グラフ要約（対 w=1 total 倍率、3 回の最小値）}
  \label{tab:shm-summary}
  \begin{tabular}{lrrrr}
    \toprule
    & \multicolumn{2}{c}{fork} & \multicolumn{2}{c}{shm} \\
    \cmidrule(lr){2-3}\cmidrule(lr){4-5}
    グラフ & w=4 & w=8 & w=4 & w=8 \\
    \midrule
    木 $n{=}20$k          & 1.97 & 1.88 & 1.22 & 0.97 \\
    2 次元格子 $100\times100$ & 2.01 & 1.92 & 1.18 & 0.86 \\
    疎ER $n{=}20$k        & 2.09 & 2.28 & 1.47 & 1.17 \\
    BA hubs $n{=}20$k     & 2.35 & 2.61 & 1.58 & 1.26 \\
    密ER $n{=}120$        & 1.08 & 1.07 & 1.07 & 0.99 \\
    Moon–Moser $k{=}11$   & 0.97 & 0.95 & 0.96 & 0.85 \\
    ca-AstroPh            & 3.09 & 3.91 & 2.35 & 2.20 \\
    ca-HepPh              & 2.76 & 2.86 & 2.23 & 1.80 \\
    soc-sign-epinions     & 2.65 & 3.05 & 2.78 & 3.24 \\
    \bottomrule
  \end{tabular}
\end{table}

\subsection{fork は startup をほぼ消した}

全 9 グラフで fork の startup は 0.002〜0.031 秒に収まった。
spawn では worker 数に比例して伸びていた固定費（soc で w=8 の 1.273 秒）が、fork では w=8 でも 0.031 秒である。
この結果、spawn では並列化が全敗だった小さい 6 グラフのうち 4 つ（木・格子・疎ER・BA）が w=4 で約 2 倍の得に転じた。
並列化の損益分岐点は「compute 数秒」から「compute 数十ミリ秒」まで下がったことになる。
残る密ER と Moon–Moser が改善しないのは startup のせいではなく、compute 自体が縮まないためである。
単一の重い部分木が支配的で、バッチ並列の粒度では割れない。

\subsection{shm は転送しか消せない}

shm の startup は spawn より小さいが fork より一桁大きい（soc w=1 で 0.413 秒、小グラフで 0.04〜0.16 秒）。
残るのは親側の CSR パック、spawn の interpreter 起動、worker 側の set 再構築である。
この実験は spawn の startup の内訳を分離したことになる。
転送を消すだけでは足りず、構造構築と interpreter 起動が残る限り、小グラフの損益分岐点は大きくは動かない。
小さい 6 グラフでの得は最大 1.58 倍（BA w=4）にとどまり、fork に全面的に劣る。

\subsection{shm の compute が fork より速い}

予想外の観測として、shm の compute が fork より一貫して速かった。
soc w=1 で 85.2 秒対 96.2 秒（11\%）、ca-AstroPh w=1 で 0.703 秒対 0.841 秒である。
再帰カーネルは同一なので、手法のコード差では説明できない。

有力な機構は fork の copy-on-write の遅延コストである。
Python はオブジェクトを読むだけでも参照カウントを書き換えるため、fork で共有したページは compute 中に次々とコピーされる。
つまり fork の startup の安さの一部は、コピーを compute に付け替えたものと読める。
一方 shm の worker は自分で作ったコンパクトな set を読むのでこのコストがない。
熱ドリフトなら後に走った shm が遅くなるはずで方向が逆なので、環境説より CoW 遅延コスト説が有力である。
ただし切り分け（例：\texttt{gc.freeze()} を使った fork の再測定）は未検証である。

\subsection{到達点}

soc の w=8 total は spawn 37.8 秒（前回測定）、fork 31.7 秒、shm 26.5 秒と縮んだ。
ただし spawn の値は別セッションの測定で、絶対値は環境に極めて敏感（Chrome 開閉で 26\%）なので、厳密な手法間比較として読んでよいのは同一セッションで隣接して測った fork 対 shm のみである。

結論として、startup 固定費を消す目的には fork が最小の変更で最大の効果を与える。
shm は startup 削減では fork に劣るが、compute を含めた total では soc で最速だった。
両者の両立（fork と共有配列の併用、あるいは \texttt{gc.freeze}）は将来課題である。
